\section{为什么“形成过程”对干预重要}

一旦把框架理解为涌现出来的历史性结构，干预设计就会发生变化。
核心问题不再只是“我现在如何抵抗坏状态？”，
而会变成“哪些重复的局部耦合仍在生产坏状态，以及怎样重设计循环，才能在后续制造出不同的宏观模式？”

这为干预提供了一个更严格的理由：
\begin{enumerate}[nosep]
  \item 识别重复的线索--动作--奖赏循环；
  \item 识别究竟是哪一个尺度在承载该机制；
  \item 识别该机制究竟是刚性的、亚稳态的，还是脆弱的；
  \item 重设计重复条件，而不是只依赖“关键时刻的英雄式抵抗”。
\end{enumerate}

这就是涌现语言的业务价值。
它把手稿从一种“抵抗坏状态”的理论，转变成一种“制造更好状态”的理论。